\subsection{Fallacy detection in NLP}
\label{sec:bg-fallacy}

The Reddit informal-fallacy dataset that the Touch\'{e} 2026 lab redistributes was introduced by Sahai et al.~\cite{sahai2021invisiblewall}, who labelled a long-tail eight-way taxonomy over comments mined from r/changemyview and adjacent subreddits and reported the first transformer baselines on the resulting corpus. Adjacent benchmarks treat closely related phenomena with similar fine-grained label structures. The Argument Reasoning Comprehension task introduced multi-class warrant identification on user-written argument pairs and surfaced how brittle large language models can be on argumentation tasks that require implicit-premise reasoning~\cite{habernal2018arc}. Da San Martino et al.~\cite{dasanmartino2019propaganda} catalogued eighteen propaganda techniques in news articles. Several of those techniques (\emph{appeal to authority}, \emph{appeal to fear}, \emph{black-and-white fallacy}) overlap directly with the ST2 inventory and indicate that the propaganda-detection literature is a useful cross-domain reference even though the source genre differs. Goffredo et al.~\cite{goffredo2022political} classified fallacies in political-debate transcripts; their taxonomy differs from the ST2 inventory, but their supervised setting is a close methodological neighbour for the classifier we build.

Two more recent papers anchor the picture. Alhindi et al.~\cite{alhindi2022multitask} reported strong multitask instruction-prompted results across a suite of fallacy-recognition datasets, demonstrating that multitask supervision over related fallacy datasets transfers usefully to eight-way tasks of this kind. MAFALDA~\cite{helwe2024mafalda} is the most-cited 2024 unified fallacy benchmark, we do not retrain on MAFALDA, but our ST2 label inventory is recognisable in its taxonomy and the MAFALDA error-analysis findings on rare-class behaviour echo what we report for the ST3 \textsf{pe} class in \S\ref{sec:disc-pe}. Pan et al.~\cite{pan2024zeroshotfallacy} provide the published precedent for evaluating LLMs as zero-shot fallacy classifiers. They show that zero-shot LLM performance is real but uneven across fallacy types, which motivates the LLM track we describe in \S\ref{sec:llm-detour} as worth running rather than dismissing.

\subsection{Argumentation schemes}
\label{sec:bg-schemes}

The four ST3 labels are not ad hoc, but neither do they map one-to-one onto individual Walton schemes. The goal (practical / epistemic) $\times$ basis (internal / external) crossing is the lab's operationalisation~\cite{kiesel:2026c} of the argumentation-profile framing of Macagno~\cite{macagno2022argumentation}, itself rooted in the broader Walton--Reed--Macagno scheme tradition~\cite{walton2008schemes}. The crossing summarises, at a coarser grain than a full scheme catalogue, whether the speaker is arguing about what to do or what to believe, and whether the warrant comes from the speaker's own resources or from an external source. NLP work that classifies arguments by Walton scheme is comparatively thin: Feng and Hirst~\cite{feng2011scheme} gave the canonical treatment, assigning arguments to five common schemes, while Jo et al.~\cite{jo2021schemes} classify argumentative \emph{relations} using scheme-based logical mechanisms rather than scheme labels. Neither targets the ST3 goal/basis label set directly, so a system designer cannot transfer a published scheme-classification recipe onto this task.

The Touch\'{e} lab itself has progressively shifted over recent editions from argument retrieval~\cite{bondarenko2022touche,bondarenko2023touche} to full argumentation systems~\cite{kiesel2024touche,kiesel2025touche}, with the 2026 edition the first to feature fallacy detection as a dedicated task. The encoder-vs-LLM trade-off for small-corpus fallacy classification is not, to our reading, settled by the published literature: encoder fine-tuning, LoRA-adapted mid-size LLMs, and zero/few-shot prompted frontier models have each been reported as competitive on neighbouring tasks, with no clean recipe-matched comparison at the corpus scale of the present shared task. That open question is the gap our contribution C3 addresses, and it motivates the systematic backbone comparison reported in \S\ref{sec:results}.
